\subsection{Búsqueda Binaria por Rango}\label{sec:binario-rango}
La búsqueda binaria por rango es una variante de la búsqueda binaria que permite determinar el intervalo de posiciones en el que aparece un valor dentro de un arreglo ordenado, o incluso un segmento de elementos cuyos valores se encuentran dentro de un rango establecido.

En este proyecto se implementaron dos variantes complementarias de esta idea:
\begin{enumerate}
    \item Búsqueda por valor exacto: recibe un único valor y determina la primera y última posición en la que aparece dentro del arreglo.
    \item Búsqueda por intervalo: recibe un valor mínimo y un valor máximo, y determina el bloque de elementos cuyos valores están entre ambos límites.
\end{enumerate}
Esta distinción resulta útil en el contexto del sistema de deportistas, ya que múltiples jugadores pueden compartir el mismo puntaje. Por eso mismo, no basta con solo encontrar una sola coincidencia, ya que, en muchos casos interesa recuperar todas las concurrencias asociadas a un valor exacto, o bien todos los deportistas cuyo puntaje pertenezcan a un intervalo definido.

\subsubsection{Búsqueda por valor exacto: \texttt{first} y \texttt{last}}
La primera variante consiste en localizar todas las ocurrencias de un valor exacto dentro de un arreglo ordenado. Para ello se implementan dos funciones auxiliares:
\begin{itemize}
    \item \texttt{first}: encuentra la primera posición en la que aparece el valor buscado.
    \item \texttt{last}: encuentra la última posición del valor buscado.
\end{itemize}

La idea consiste en ejecutar dos búsquedas binarias modificadas. En ambos casos, cuando se encuentra una coincidencia, el algoritmo no termina inmediatamente, sino que continúa explorando hacia la izquierda o hacia la derecha para hallar el extremo correspondiente. De esta manera \texttt{first} sigue buscando hacia la izquierda mientras existan coincidencias anteriores, y \texttt{last} sigue buscando hacia la derecha mientras existan coincidencias posteriores. Una vez obtenido ambos índices, el rango completo de ocurrencias queda determinado por el intervalo \texttt{[first, last]}.

\subsubsection{Búsqueda por intervalo: \texttt{lower\_bound} y \texttt{upper\_bound}}
La segunda variante corresponde a la búsqueda de todos los elementos cuyos valores se encuentran dentro de un rango \texttt{[min, max]}. En este caso no se busca un valor exacto, sino los límites del segmento válido dentro del arreglo ordenado.

Para ello se implementan las siguientes dos funciones auxiliares:
\begin{itemize}
    \item \texttt{lower\_bound}: determina la primera posición cuyo valor es mayor o igual que el límite inferior.
    \item \texttt{upper\_bound}: determina la última posición cuyo valor es menor o igual que el límite superior.
\end{itemize}
Estas funciones permiten identificar directamente el subarreglo que satisface la condición del rango, evitando recorrer elementos innecesarios. Si el índice retornado por \texttt{lower\_bound} supera el retornado por \texttt{upper\_bound}, entonces el rango buscado no contiene elementos válidos.

\subsubsection{Idea general en pseudocódigo}
A continuación se presentan pseudocódigos que describen de manera abstracta el comportamiento de estas variantes de búsqueda binaria por rango. Para observar su implementación exacta, puede consultarse el código fuente disponible en el repositorio.

En la búsqueda por valor exacto, el proceso es:
\begin{algorithm}
    \caption{Rango de apariciones $\triangleright$ $V[0..n-1]$}
    \begin{algorithmic}[1]
        \Procedure{SearchRange}{$V, x$}
        
        \State $first \gets \Call{FirstOcurrence}{V, x}$

        \If{$first = -1$}
            \State \Return $\emptyset$
        \EndIf

        \State $last \gets \Call{LastOcurrence}{V, x}$
        \State \Return $[first, last]$

        \EndProcedure
    \end{algorithmic}
\end{algorithm}

En la búsqueda por intervalo, el procedimiento es:
\begin{algorithm}
    \caption{Rango de valores entre mínimo y máximo $\triangleright$ $V[0..n-1]$}
    \begin{algorithmic}[1]
        \Procedure{SearchRangeBetween}{$V, min, max$}
        
        \State $first \gets \Call{LowerBound}{V, min}$
        \State $last \gets \Call{UpperBound}{V, max}$

        \If{$first > last$}
            \State \Return $\emptyset$
        \EndIf

        \State \Return $[first, last]$

        \EndProcedure
    \end{algorithmic}
\end{algorithm}

En ambos casos, el resultado final no es un único índice, sino un intervalo de posiciones que representa todas las coincidencias con el valor o segmento solicitado.

\subsubsection{Análisis de complejidad}
Cada una de las funciones auxiliares (\texttt{first}, \texttt{last}, \texttt{lower\_bound}, \texttt{upper\_bound}) realiza una búsqueda binaria sobre un arreglo ordenado, por lo que su complejidad es: $$O(log\text{ }n)$$ Como cada consulta requiere dos búsquedas de este tipo, el costo total sigue siendo: $$O(log\text{ }n)$$ ya que la suma de dos términos logarítmicos mantiene el mismo orden de crecimiento.

Por tanto, ambas variantes:
\begin{itemize}
    \item Mejor caso: $O(log\text{ }n)$
    \item Caso promedio: $O(log\text{ }n)$
    \item Peor caso: $O(log\text{ }n)$
\end{itemize}

Esto representa una mejora importante respecto de una solución lineal, que requeriría recorrer potencialmente todo el arreglo para identificar el conjunto de elementos buscados.

\subsubsection{Comentario general}
La búsqueda binaria por rango aumenta la útilidad de la búsqueda binaria clásica, permitiendo recuperar conjuntos de elementos en lugar de una sola coincidiencia. En este proyecto, su implementación en dos variantes diferenciadas, una basada en \texttt{first} y \texttt{last} para valores exactos, y otra basada en \texttt{lower\_bound} y \texttt{upper\_bound} para intervalos, permite resolver consultas más expresivas sobre los puntajes de los deportistas sin abandonar la eficiencia logarítmica del enfoque divide y vencerás.
